\pdfinfoomitdate=1
\pdftrailerid{}
\pdfsuppressptexinfo=-1
\documentclass[11pt]{article}
\usepackage[margin=1.02in]{geometry}
\usepackage{amsmath,amssymb,amsthm,mathtools,booktabs,array,microtype}
\usepackage[hidelinks]{hyperref}
\usepackage[nameinlink,capitalise]{cleveref}
\usepackage{enumitem}

\newtheorem{theorem}{Theorem}[section]
\newtheorem{proposition}[theorem]{Proposition}
\newtheorem{lemma}[theorem]{Lemma}
\newtheorem{corollary}[theorem]{Corollary}
\theoremstyle{remark}
\newtheorem{remark}[theorem]{Remark}

\DeclareMathOperator{\Hess}{Hess}
\DeclareMathOperator{\per}{per}
\DeclareMathOperator{\Sym}{Sym}
\DeclareMathOperator{\GL}{GL}
\DeclareMathOperator{\SL}{SL}
\newcommand{\C}{\mathbb C}
\newcommand{\Q}{\mathbb Q}
\newcommand{\Dcal}{\mathcal D}
\newcolumntype{L}[1]{>{\raggedright\arraybackslash}p{#1}}

\title{A Uniform Line Formula for the Hessian Determinant of the Permanent\\
and Equal-Size Hessian-Orbit Incomparability}
\author{Sasha Lempers}
\date{July 16, 2026}

\begin{document}
\maketitle

\begin{abstract}
Let $V=M_n(\C)$ and let $\GL(V)$ act on forms by linear change of variables.  For the relative Hessian covariant
\[
 \Dcal(f)=\det\Hess(f),
\]
we prove the exact restriction formula
\[
 \Dcal(\per_n)\bigl(J_n+(t-1)E_{11}\bigr)
 =\frac{((n-2)!)^{n^2}(n-1)^{2n}}{(n-2)^{(n-2)^2}}
   (t+n-3)^{(n-2)^2}
\]
for every $n\geq3$.  Combined with the identity
\[
 \Dcal(\det_n)
 =(n-1)(-1)^{(n-1)(n+2)/2}\det_n^{\,n(n-2)},
\]
a closed power cone, and a verified GIT polystability criterion, this yields both equal-size noncontainments
\[
 \Dcal(\det_n)\notin\overline{\GL(V)\cdot\Dcal(\per_n)},
 \qquad
 \Dcal(\per_n)\notin\overline{\GL(V)\cdot\Dcal(\det_n)}.
\]
A relative-covariant transfer principle then recovers the corresponding equal-size determinant/permanent noncontainments.  That latter corollary is not claimed as new: it is already an immediate inference from published stabilizer and polystability results.  No priority claim is made for the line formula or the Hessian-orbit theorem pending a broader specialist literature review.

The accompanying repository reconstructs the size-three identities with three distinct computational engines.  It confirms $\Dcal(\det_3)=-2\det_3^3$, $\Dcal(\per_3)(J_3+(t-1)E_{11})=64t$, squarefreeness of $\Dcal(\per_3)$, and the two complete ordinary catalecticant rank vectors.  The result is equal-size and unpadded.  It gives no determinantal-complexity lower bound and no separation of $\mathrm{VP}$ from $\mathrm{VNP}$; the full ambient Hessian determinant vanishes on the standard padded permanent whenever unused variables are present.
\end{abstract}

\section{Statement, conventions, and claims}
Set
\[
 V=M_n(\C),\qquad N=n^2,\qquad q=N(n-2),\qquad r=n(n-2).
\]
For $f\in\Sym^nV^*$, define $\Dcal(f)=\det\Hess(f)\in\Sym^qV^*$.  The main result concerns the images under $\Dcal$, not the padded permanent orbit closure used in the standard GCT formulation.

\begin{theorem}[Equal-size Hessian-orbit incomparability]\label{thm:main}
For every $n\geq3$,
\[
 \Dcal(\det_n)\notin\overline{\GL(V)\cdot\Dcal(\per_n)}
 \quad\text{and}\quad
 \Dcal(\per_n)\notin\overline{\GL(V)\cdot\Dcal(\det_n)}.
\]
Consequently the two Hessian-image orbit closures are incomparable.
\end{theorem}

The proof is entirely symbolic.  Finite computations are recorded separately and are not substituted for any uniform argument.

\section{Relative covariance and transfer of containment}\label{sec:transfer}

\begin{lemma}[Hessian covariance]\label{lem:covariance}
For $g\in\GL(V)$, under $(g\cdot f)(x)=f(g^{-1}x)$,
\[
 \Dcal(g\cdot f)=\det(g)^{-2}\,g\cdot\Dcal(f).
\]
\end{lemma}
\begin{proof}
The chain rule gives
\[
 \Hess(g\cdot f)(x)=g^{-T}\Hess(f)(g^{-1}x)g^{-1}.
\]
Taking determinants proves the formula.
\end{proof}

\begin{lemma}[Twisted and ordinary nonzero orbits]\label{lem:twisted}
Let $0\neq u\in\Sym^qV^*$, where $q>0$.  Then
\[
 \{\det(g)^{-2}g\cdot u:g\in\GL(V)\}=\GL(V)\cdot u.
\]
\end{lemma}
\begin{proof}
For one inclusion, fix $g$ and choose $\lambda\in\C^*$ with $\lambda^{-q}=\det(g)^{-2}$.  Since $(\lambda I)\cdot u=\lambda^{-q}u$,
\[
 \det(g)^{-2}g\cdot u=(\lambda I)g\cdot u.
\]
Conversely, choose $\lambda\in\C^*$ satisfying $\lambda^{q+2N}=\det(g)^{-2}$.  Then
\[
 \det(\lambda g)^{-2}(\lambda g)\cdot u=g\cdot u.
\]
All required roots exist over $\C$.
\end{proof}

\begin{proposition}[Relative-covariant transfer principle]\label{prop:transfer}
For $f,p\in\Sym^nV^*$,
\[
 f\in\overline{\GL(V)\cdot p}
 \quad\Longrightarrow\quad
 \Dcal(f)\in\overline{\GL(V)\cdot\Dcal(p)}.
\]
\end{proposition}
\begin{proof}
The map $\Dcal$ is polynomial in the coefficients, hence Zariski-continuous.  If $\Dcal(p)=0$, covariance shows that $\Dcal$ vanishes on $\GL(V)\cdot p$, and continuity therefore gives $\Dcal(f)=0$; the conclusion is immediate.  Assume henceforth that $\Dcal(p)\neq0$.  Applying $\Dcal$ to the assumed containment, using \cref{lem:covariance}, and then using \cref{lem:twisted}, identifies the closure of the twisted image orbit with the ordinary $\GL(V)$-orbit closure of $\Dcal(p)$.
\end{proof}

\section{The determinant identity and the permanent line formula}\label{sec:formulas}

\begin{theorem}[Hessian determinant of the determinant]\label{thm:detformula}
For every $n\geq2$,
\[
 \Dcal(\det_n)=c_n\det_n^{\,r},
 \qquad
 c_n=(n-1)(-1)^{(n-1)(n+2)/2}.
\]
\end{theorem}
\begin{proof}
For $X\in\GL_n$ and $U,W\in M_n$, Jacobi's formula yields
\[
 d^2(\det)_X(U,W)
 =\det(X)\bigl(\operatorname{tr}(X^{-1}U)\operatorname{tr}(X^{-1}W)
 -\operatorname{tr}(X^{-1}UX^{-1}W)\bigr).
\]
At $I$, the Hessian bilinear form is
\[
 B(U,W)=\operatorname{tr}(U)\operatorname{tr}(W)-\operatorname{tr}(UW).
\]
At $X$, the Hessian is $\det(X)$ times the pullback of $B$ under left multiplication by $X^{-1}$.  Since $\det(L_{X^{-1}})=\det(X)^{-n}$,
\[
 \det\Hess(\det_n)(X)=\det(B)\det(X)^{n^2-2n}.
\]
On diagonal matrices, $B$ has matrix $J_n-I_n$, whose determinant is $(n-1)(-1)^{n-1}$.  Each subspace $\langle E_{ij},E_{ji}\rangle$, $i<j$, contributes the block
\[
 \begin{pmatrix}0&-1\\-1&0\end{pmatrix}
\]
of determinant $-1$.  There are $n(n-1)/2$ such blocks, giving the stated sign.  Equality on the dense set $\GL_n$ extends polynomially to all $M_n$.
\end{proof}

Let $X(t)=J_n+(t-1)E_{11}$, and set
\[
 A=J_n-I_n,
 \qquad
 C=0\oplus(J_{n-1}-I_{n-1}).
\]

\begin{theorem}[Uniform permanent line formula]\label{thm:line}
For every $n\geq3$,
\[
 \Dcal(\per_n)(X(t))=K_n(t+n-3)^{(n-2)^2},
\]
where
\[
 K_n=\frac{((n-2)!)^{n^2}(n-1)^{2n}}{(n-2)^{(n-2)^2}}\neq0.
\]
\end{theorem}
\begin{proof}
A mixed second derivative of $\per_n$ is zero when the two differentiated entries share a row or a column; otherwise it equals the permanent of the complementary $(n-2)\times(n-2)$ minor.  Evaluating at $X(t)$ gives
\[
 \Hess(\per_n)(X(t))
 =(n-2)!A\otimes A+(t-1)(n-3)!C\otimes C.
\]
Put $B=A^{-1}C$.  We now prove the spectrum rather than merely assert it.  Let
\[
 W=\{x\in\C^n:x_1=0,\ \sum_{i=2}^n x_i=0\},
 \qquad u=(0,1,\ldots,1)^T.
\]
For $x\in W$, one has $Ax=Cx=-x$, and therefore $Bx=x$.  Moreover $Be_1=0$.  Since $Cu=(n-2)u$ and $A^{-1}u=e_1$,
\[
 Bu=(n-2)e_1.
\]
Thus $B$ is the identity on the $(n-2)$-dimensional space $W$ and has a nilpotent $2\times2$ block on $\langle e_1,u\rangle$.  Its spectrum is consequently $1$ with algebraic multiplicity $n-2$ and $0$ with algebraic multiplicity $2$.

Factoring the first Kronecker term gives
\[
 \det\Hess(\per_n)(X(t))
 =((n-2)!)^{n^2}\det(A\otimes A)
 \det\left(I+\frac{t-1}{n-2}B\otimes B\right).
\]
Now $\det(A\otimes A)=\det(A)^{2n}=(n-1)^{2n}$.  The operator $B\otimes B$ has eigenvalue $1$ with algebraic multiplicity $(n-2)^2$, and all its other eigenvalues are zero.  Hence
\[
 \det\left(I+\frac{t-1}{n-2}B\otimes B\right)
 =\left(\frac{t+n-3}{n-2}\right)^{(n-2)^2},
\]
which is the claimed formula.
\end{proof}

\begin{corollary}[Nonpower property]\label{cor:nonpower}
The form $\Dcal(\per_n)$ is not a nonzero scalar multiple of an $r$-th power, where $r=n(n-2)$.
\end{corollary}
\begin{proof}
A nonzero scalar multiple of an $r$-th power has, on every line where it is not identically zero, root multiplicities divisible by $r$.  The restriction in \cref{thm:line} has one root of multiplicity $(n-2)^2$, and
\[
 0<(n-2)^2<n(n-2)=r.
\]
\end{proof}

\begin{lemma}[Closed power cone]\label{lem:powercone}
The set
\[
 \mathcal P_r=\{h^r:h\in\Sym^nV^*\}\subseteq\Sym^qV^*
\]
is a closed $\GL(V)$-stable affine cone.
\end{lemma}
\begin{proof}
Its projectivization is the image of the Veronese embedding $[h]\mapsto[h^r]$, hence is closed.  Taking the affine cone preserves closedness.  Stability follows from $g\cdot h^r=(g\cdot h)^r$.
\end{proof}

\section{Polystability of the Hessian forms}\label{sec:polystability}
We use the support-cone criterion of B\"urgisser and Ikenmeyer: a form is $\SL$-polystable when it has a reductive stabilizer subgroup with diagonal centralizer and when the cone generated by its monomial support contains the all-ones vector \cite[Proposition~2.8]{BI}.

\begin{proposition}[Polystability]\label{prop:polystability}
For every $n\geq3$, both $\Dcal(\det_n)$ and $\Dcal(\per_n)$ have closed $\SL(V)$-orbits.
\end{proposition}
\begin{proof}
Let $R\subseteq\SL(V)$ be the row-column torus acting on $E_{ij}$ through the character $(a,b)\mapsto a_i b_j$, with $\prod_i a_i=\prod_j b_j=1$.  It fixes $\det_n$ and $\per_n$, and therefore fixes their Hessian determinants by \cref{lem:covariance}.

The characters are pairwise distinct.  Indeed, for $(i,j)\neq(k,l)$, if $i\neq k$ vary the two relevant row parameters by $a_i=s$, $a_k=s^{-1}$ and leave the column parameters equal to one.  If $i=k$, then $j\neq l$ and the analogous variation of $b_j,b_l$ separates the characters.  Therefore the $R$-weight spaces in $V$ are exactly the coordinate lines $\C E_{ij}$, so the centralizer of $R$ is contained in the diagonal torus.

The two Hessian forms are nonzero by \cref{thm:detformula,thm:line}.  Row and column permutations preserve each monomial support up to an overall scalar and act transitively on the $N=n^2$ variables.  Averaging any support exponent vector of total degree $q$ over this finite group yields
\[
 (q/N)(1,\ldots,1)=(n-2)(1,\ldots,1).
\]
Thus the support cone contains $(1,\ldots,1)$, and the cited criterion applies.
\end{proof}

\begin{lemma}[Polystable degeneration]\label{lem:closedorbit}
Let $0\neq u,v\in\Sym^qV^*$ be $\SL(V)$-polystable.  If $u\in\overline{\GL(V)\cdot v}$, then $u\in\GL(V)\cdot v$.
\end{lemma}
\begin{proof}
A point in the Zariski closure of a complex orbit is a Euclidean limit of orbit points.  Hence, after absorbing the central action into a scalar, there are $a_k\in\C^*$ and $s_k\in\SL(V)$ with
\[
 a_k s_k\cdot v\longrightarrow u.
\]
Because $u$ is nonzero and polystable, some homogeneous $\SL(V)$-invariant polynomial $P$ of degree $e>0$ satisfies $P(u)\neq0$.  Then
\[
 a_k^eP(v)\longrightarrow P(u),
\]
so $P(v)\neq0$ and $a_k^e$ converges to the nonzero number $P(u)/P(v)$.  Therefore $(a_k)$ is bounded and bounded away from zero.  After taking a subsequence, $a_k$ converges to a nonzero $e$-th root $a$ of that limit.  It follows that
\[
 s_k\cdot v\longrightarrow a^{-1}u.
\]
Since $\SL(V)\cdot v$ is closed, $a^{-1}u\in\SL(V)\cdot v$, and hence $u\in\GL(V)\cdot v$.
\end{proof}

\section{Proof of the Hessian-orbit theorem}
\begin{proof}[Proof of \cref{thm:main}]
By \cref{thm:detformula}, $\Dcal(\det_n)$ belongs to the closed $\GL(V)$-stable cone $\mathcal P_r$.  Consequently
\[
 \overline{\GL(V)\cdot\Dcal(\det_n)}\subseteq\mathcal P_r.
\]
By \cref{cor:nonpower}, $\Dcal(\per_n)\notin\mathcal P_r$.  This proves
\[
 \Dcal(\per_n)\notin\overline{\GL(V)\cdot\Dcal(\det_n)}.
\]

The two Hessian forms are not orbit-equivalent, since one is an $r$-th power up to scalar and the other is not.  They are both polystable by \cref{prop:polystability}.  If
\[
 \Dcal(\det_n)\in\overline{\GL(V)\cdot\Dcal(\per_n)},
\]
then \cref{lem:closedorbit} would force orbit equivalence, a contradiction.  The reverse noncontainment follows.
\end{proof}

\section{The original equal-size corollary and its prior status}

\begin{corollary}[Equal-size determinant/permanent incomparability]\label{cor:original}
For every $n\geq3$,
\[
 \det_n\notin\overline{\GL(V)\cdot\per_n}
 \quad\text{and}\quad
 \per_n\notin\overline{\GL(V)\cdot\det_n}.
\]
\end{corollary}
\begin{proof}
Apply the contrapositive of \cref{prop:transfer} to the two noncontainments in \cref{thm:main}.
\end{proof}

\begin{remark}[No novelty claim for \cref{cor:original}]
B\"urgisser and Ikenmeyer record the classical stabilizers of $\det_n$ and $\per_n$ and prove their polystability \cite[Theorems~2.5--2.6 and Section~2.2]{BI}.  The identity components of the two stabilizers have dimensions $2(n^2-1)$ and $2(n-1)$, respectively, so the forms are not orbit-equivalent for $n\geq3$.  The same closed-orbit argument as \cref{lem:closedorbit} then already gives \cref{cor:original}.  Thus the corollary is included only for logical completeness and is not presented as a new theorem.
\end{remark}

\section{Padding barrier and exact relation to GCT}\label{sec:padding}
For $m>n$, the standard homogeneous comparison uses
\[
 P_{n,m}=\ell^{m-n}\per_n
\]
as a degree-$m$ form in the $m^2$-variable determinant ambient space.

\begin{proposition}[Ambient Hessian padding barrier]\label{prop:padding}
For $m>n$,
\[
 \det\Hess_{\C^{m^2}}(P_{n,m})=0.
\]
\end{proposition}
\begin{proof}
The padded form depends on at most $n^2+1$ variables.  Since $m>n$ implies $m^2>n^2+1$, at least one ambient direction is unused.  The corresponding Hessian row and column are zero.
\end{proof}

The GCT determinant-versus-padded-permanent formulation compares orbit closures in this larger ambient space \cite{MS1,MS2,BI}.  Occurrence-obstruction no-go theorems and shifted-partial-derivative limitations apply to that padded regime \cite{BIP,ELSW}.  The present Hessian theorem neither contradicts nor bypasses them.  In particular, it establishes none of the following:
\begin{itemize}[leftmargin=2em]
\item a superpolynomial lower bound for determinantal or border determinantal complexity;
\item a noncontainment for the padded permanent orbit closure;
\item a separation of $\mathrm{VP}$ from $\mathrm{VNP}$.
\end{itemize}

\section{Exact size-three reconstruction}
The repository supplies three distinct catalecticant engines: a primary exact SymPy implementation, a clean-room exponent-dictionary implementation with an independent polynomial determinant, and a pure-Python modular implementation with its own sparse elimination.  Their script hashes are distinct.  They agree on the canonical polynomial hashes, all direct integer matrix hashes for orders $0$ through $4$, and all rank vectors.

The exact reconstructed identities are
\[
 \Dcal(\det_3)=-2\det_3^3,
 \qquad
 \Dcal(\per_3)(J_3+(t-1)E_{11})=64t.
\]
The form $\Dcal(\per_3)$ has degree $9$ and $55$ monomials.  Exact factorization over $\Q$ returns one degree-$9$ factor of multiplicity one; no absolute irreducibility claim is made.  The exact gcd over $\Q$ of its nine first partial derivatives is $1$, which proves squarefreeness after extension to $\C$.

The complete ordinary catalecticant rank vectors are
\begin{align*}
 (\operatorname{rank}C_a(\Dcal(\det_3)))_{a=0}^{9}
   &=(1,9,45,165,270,270,165,45,9,1),\\
 (\operatorname{rank}C_a(\Dcal(\per_3)))_{a=0}^{9}
   &=(1,9,45,165,414,414,165,45,9,1).
\end{align*}
Therefore every closed condition obtained solely by imposing upper bounds on ordinary catalecticant ranks and containing the permanent-Hessian orbit closure also contains $\Dcal(\det_3)$.  This is a statement about these rank loci only; it is not an orbit-containment theorem and says nothing by itself about Young or Koszul flattenings, syzygies, or multiplicity obstructions.

\section{Reproducibility contract}
The canonical serialization uses the fixed variable order
\[
 x_{11},x_{12},x_{13},x_{21},x_{22},x_{23},x_{31},x_{32},x_{33},
\]
Exponent vectors are ordered lexicographically in descending order.  Rational coefficients are serialized as numerator--denominator pairs in compact UTF-8 JSON.  Catalecticant bases are weak compositions with the first coordinate increasing recursively.  Certificates contain no elapsed times, absolute paths, host names, or timestamps.

From the repository root, the complete check is
\begin{verbatim}
./reproduce_all.sh
\end{verbatim}
It rebuilds the certificates, compares two fresh runs byte for byte, executes the consistency tests, and verifies \texttt{certificates/SHA256SUMS.txt}.  The term ``computational certificate'' here means an exact, deterministic, independently reconstructed calculation; it does not mean a kernel-checked formal proof.

\section{Claim and priority audit}
\begin{center}
\small
\begin{tabular}{@{}L{.43\linewidth}L{.20\linewidth}L{.28\linewidth}@{}}
\toprule
Claim & Status & Exact scope \\
\midrule
Determinant Hessian identity & proved & all $n\geq2$ \\
Permanent line formula & proved & all $n\geq3$ \\
Nonpower property & proved & all $n\geq3$ \\
Polystability of Hessian forms & proved using published criterion & all $n\geq3$ \\
Hessian-orbit incomparability & proved & equal-size, unpadded \\
Original equal-size incomparability & literature-level inference; reproved here & not claimed new \\
Size-three ranks and squarefreeness & exact reproducible computation & three distinct engines for ranks \\
Priority of line formula/theorem & unknown & no ``first'' or ``new'' claim \\
Padded lower bound or $\mathrm{VP}\neq\mathrm{VNP}$ & not obtained & ambient Hessian vanishes \\
\bottomrule
\end{tabular}
\end{center}

\begin{thebibliography}{99}
\bibitem{MS1}
K.~D. Mulmuley and M. Sohoni,
\emph{Geometric Complexity Theory I: An Approach to the P vs. NP and Related Problems},
SIAM J. Comput. 31 (2001), 496--526.
DOI: \href{https://doi.org/10.1137/S009753970038715X}{10.1137/S009753970038715X}.

\bibitem{MS2}
K.~D. Mulmuley and M. Sohoni,
\emph{Geometric Complexity Theory II: Towards Explicit Obstructions for Embeddings among Class Varieties},
SIAM J. Comput. 38 (2008), 1175--1206.
DOI: \href{https://doi.org/10.1137/080718115}{10.1137/080718115}.

\bibitem{BI}
P. B\"urgisser and C. Ikenmeyer,
\emph{Fundamental invariants of orbit closures},
J. Algebra 477 (2017), 390--434.
DOI: \href{https://doi.org/10.1016/j.jalgebra.2016.12.035}{10.1016/j.jalgebra.2016.12.035}; arXiv:1511.02927.

\bibitem{BIP}
P. B\"urgisser, C. Ikenmeyer, and G. Panova,
\emph{No occurrence obstructions in geometric complexity theory},
J. Amer. Math. Soc. 32 (2019), 163--193.
DOI: \href{https://doi.org/10.1090/jams/908}{10.1090/jams/908}; arXiv:1604.06431.

\bibitem{ELSW}
K. Efremenko, J.~M. Landsberg, H. Schenck, and J. Weyman,
\emph{The method of shifted partial derivatives cannot separate the permanent from the determinant},
Math. Comp. 87 (2018), 2037--2045.
DOI: \href{https://doi.org/10.1090/mcom/3284}{10.1090/mcom/3284}; arXiv:1504.05171.

\bibitem{Fox}
D.~J.~F. Fox,
\emph{Functions dividing their Hessian determinants and affine spheres},
Asian J. Math. 20 (2016), 503--530; arXiv:1307.5394.
\end{thebibliography}

\end{document}
